% SGA 3，第 VII_B 篇，英文重建稿——第 1 轮。
% 范围：权威底本分篇第 95 页，定义 5.2.6 全文，止于第 5.3 节之前。
% 合并阅读本页码： 608.
% 再版印刷页码： 597.
% 本范围内编者注： 164.
% 权威底本： Polo--Gille Exp7B-13oct24.pdf,
% SHA-256 D303690E0195654BBE7C3DAA5D89EA23E24CD709A3575217234AC19AE8836694.
% 曾参考 Jacob Reinhold 的 CC BY 4.0 Markdown 作为比对谱系。
% 本组件没有图表；全部数学均为原生 LaTeX。

\subsubsection*{5.2.6}
\phantomsection
\label{sga3:VIIB:5.2.6}

\noindent\textbf{定义 5.2.6.}
\footnote{\label{sga3:VIIB:nde:164}编者增加了此编号，以定义 \SGARefLink{sga3:VIIB:5.3:item:a}{5.3(a)} 中使用的无限张量积。}
设 \((A_\lambda)_{\lambda\in\Lambda}\) 为一族预有限 \(k\)-代数，每个都带有增广
\[
  \varepsilon_\lambda:A_\lambda\longrightarrow k.
\]
\textit{（特别地，若每个 \(A_\lambda\) 都是剩余域为 \(k\) 的局部环，此条件成立。）}
定义无限完备化张量积
\[
  \mathcal A
  =
  \mathop{\widehat{\bigotimes}}\nolimits^{c}_{\lambda\in\Lambda}
  A_\lambda
\]
为 \(\operatorname{Alp}/k\) 中各
\[
  \mathcal A_F
  =
  \mathop{\widehat{\bigotimes}}\nolimits^{c}_{\lambda\in F}
  A_\lambda,
\]
的逆极限，其中 \(F\) 遍历 \(\Lambda\) 的有限子集。对 \(F'=F\cup\{\lambda\}\)，转移态射
\[
  \mathcal A_{F'}\longrightarrow\mathcal A_F
\]
为
\[
  \operatorname{id}\widehat\otimes\varepsilon_\lambda.
\]

特别地，假设 \(\Lambda=\mathbb N^*\)，并记
\[
  X_n=\operatorname{Spf}(A_n).
\]
则 \(\operatorname{Spf}(\mathcal A)\) 表示如下函子：对每个 \(C\in\operatorname{Alf}/k\)，联系 \(\prod_{n\geq1}X_n(C)\) 中的\emph{有限}序列的集合；即序列
\[
  (x_1,x_2,\ldots)\in\prod_{n\geq1}X_n(C)
\]
满足对充分大的 \(n\) 有
\[
  x_n=\varepsilon_n
\]
这里滥用记号，用 \(\varepsilon_n\) 表示
\[
  \varepsilon_n:A_n\longrightarrow k
\]
与结构态射 \(k\to C\) 的复合。

若再假设每个 \(A_n\) 为某个代数 \(k[[\omega'_n]]\) 的商，可用 \(0\) 表示在 \(\omega'_n\) 上为零的唯一态射
\[
  A_n\longrightarrow C
\]
于是，上面的描述就成为对充分大的 \(n\) 满足 \(x_n=0\) 的序列的集合。
